\section{Discussion}
\label{sec:discussion}

\subsection{Comparison Table}
\label{sec:compare-table}

\Cref{tab:scheme-compare} situates \TW{} against representative authenticated
encryption schemes along the axes that distinguish it: the underlying primitive,
whether the mode is parallelizable, the key, nonce, and tag sizes, associated-data
support, nonce-misuse resistance, full-input key commitment (\textsf{CMT-4}), and
context discoverability (CDY).
The comparison is structural; \TW{}'s own concrete bounds and the regimes in
which each term dominates are treated in \Cref{sec:bound-discussion}, and its
measured throughput against hardware-accelerated AES-128-GCM in
\Cref{sec:compare-perf}. The \textsf{CMT-4} column records full-input
commitment: \TW{} provides it within the mode, as does Ascon \cite{KSW23}, a
committing transform supplies it for an arbitrary base scheme, and the remaining
listed schemes do not provide it. The CDY column records context-discovery
resistance: \TW{} provides it within the mode (\Cref{thm:cdy}), and the
remaining entries are marked ``?'' because, to our knowledge, no
discoverability result---attack or proof---has been published for them. Context
discoverability is a recent notion \cite{MLGR23}, so its absence of entries
reflects the state of the literature rather than a negative result; the schemes
\cite{MLGR23} do break on discoverability (\textsf{GCM}, \textsf{CCM},
\textsf{EAX}, \textsf{SIV}, \textsf{OCB3}) are not among the rows above.
The table does not separately track tag-projected compactness, since that
projection strengthening has not been systematically analyzed for the comparison
schemes.

\begin{table}[t]
\centering
\footnotesize
\caption{Structural comparison of \TW{} with representative AEADs. Sizes are in
bytes; AES-GCM-SIV also admits a $32$-byte key. ``Misuse-res.'' is nonce-misuse
resistance in the sense of \cite{RS06}; the \textsf{CMT-4} column records
full-input commitment in the sense of \cite{BH22}, abstracting over its
decryption (\textsf{CMTD}) and encryption (\textsf{CMT}) forms, of which \TW{}'s
is the decryption form \textsf{CMTDs-4}; CDY is context discoverability in
the sense of \cite{MLGR23}. A dash means the entry is inherited
from the base scheme, and ``?'' that no result is established in the literature.
\TW{}'s concrete bounds appear in
\Cref{sec:bound-discussion}.}
\label{tab:scheme-compare}
\begin{tabular}{@{}llcccccc@{}}
\toprule
Scheme & Primitive & Par. & K/N/T & AD & Misuse & \textsf{CMT-4} & CDY \\
\midrule
\TW{}                        & \Keccakp   & yes & $32/32/32$ & yes & no  & yes & yes \\
Ascon-128a \cite{DEMS21}     & Ascon-$p$ ($320$ b) & no  & $16/16/16$ & yes & no  & yes & ?   \\
Xoodyak \cite{DHPVV20}       & Xoodoo ($384$ b)   & no  & $16/16/16$ & yes & no  & no  & ?   \\
AES-GCM-SIV \cite{GLL17}     & AES                & yes & $16/12/16$ & yes & yes & no  & ?   \\
nAE $+$ \textsf{CTX} \cite{CR22} & any nAE        & ---  & ---        & yes & --- & yes & ?   \\
\bottomrule
\end{tabular}
\end{table}

\subsection{Comparison with Prior Sponge AEADs}
\label{sec:compare-sponge}

\TW{} belongs to the family of single-permutation duplex AEADs initiated by
\cite{BDPV11}. The dedicated lightweight members of that family, Ascon
\cite{DEMS21} and Xoodyak \cite{DHPVV20}, target $128$-bit security on small
permutations ($320$ and $384$ bits) with a strictly serial duplex: each call
depends on the previous one, so there is no mode-level parallelism to exploit.
\TW{} instead pays for a larger permutation, \Keccakp, and a $256$-bit
capacity---the same round count as KangarooTwelve \cite{BDPVVV16}---in exchange
for a tree that processes leaf chunks in parallel and for the headroom that
yields $128$-bit security with a comfortable margin. The cost of this choice is
visible in \Cref{sec:performance}: on short, single-leaf inputs \TW{} is slower
per byte than a compact serial duplex would be, and its advantage appears only
once a message spans enough leaves to fill the vector units.

\TW{} reuses KangarooTwelve's final-node/leaf split but not its Sakura coding,
which renders that scheme's tree hashing sound: the frame bits and coded child
count of a Sakura coding make every node input unambiguously decodable, so that
any collision in the hash output entails a collision in the inner function. \TW{}
adopts no such coding and does not need one. Its tree has a single fixed
two-level shape---one root and a leaf count fixed by $\|C\|$ through
$\leaves(\cdot)$---so it never distinguishes among a family of admissible
topologies, the generality a Sakura coding exists to serve; the leaf count is
moreover pinned by absorbing $\langle n \rangle_8$ before the root tag. The
unambiguous parsing such a coding would supply is instead furnished by the
per-call domain separation of \Cref{sec:domain-sep}: the prefixes $\prfxroot$,
$\prfxleaf$ and the leaf index $\langle j \rangle_8$ fix the duplex instance, and
the seven pairwise-distinct phase suffixes of \Cref{tab:suffixes} fix the phase
and the final/non-final role of every call, so distinct transcripts flatten to
distinct sponge inputs (\Cref{lem:transcript-inj,cor:output-sep}); binding the
index into each leaf is in fact a finer separation than ordering leaves within an
aggregation alone. More to the point, the committing and discoverability
guarantees of \Cref{sec:cmtds3} do not rest on soundness of a tree coding at all:
distinct $(K, N, A)$ triples already diverge in the root initialization and
associated-data calls, before any leaf tag is aggregated (\Cref{cor:triple-sep}),
so agreement among leaf tags cannot merge distinct contexts. The leaf encoding
enters only the tag-projected bound of \Cref{thm:tag-cmtds4}, and there only
through its explicit leaf tag collision term. A Sakura coding would re-establish
parsing facts \TW{} already possesses, and would leave these bounds unchanged.

Keyak \cite{BDPVV16}, whose Motorist mode is already parallelizable, is the
closest prior duplex AEAD in spirit; \Cref{sec:related} contrasts its
wire-visible, encryptor-and-decryptor-shared degree of parallelism with \TW{}'s
run-time leaf choice.
Strobe \cite{Ham17} is a serial sponge framework aimed at whole protocols rather
than at a standalone AEAD, and like Ascon and Xoodyak it does not parallelize.
Committing security is not unique to \TW{} in this family: a dedicated analysis
of the sponge-based NIST lightweight finalists proves Ascon committing while
breaking others \cite{KSW23}, and some constructions, such as Strobe, have not
been examined in that setting at all. What distinguishes \TW{} is that it targets
the strongest $s$-way full-input notion (\textsf{CMT-4}) as a proven, designed
property of the mode, together with its preimage counterpart, context
discoverability, established without leaving the single-permutation duplex
setting; the discoverability guarantee in particular is not known to hold for
the other members of the family.

\subsection{Comparison with Committing AEAD Constructions}
\label{sec:compare-commit}

The dominant route to committing AEAD is a generic transform over a
non-committing base scheme. \cite{BH22} give the \textsf{UtC}, \textsf{RtC}, and
\textsf{HtE} transforms, which add key- or context-commitment to an arbitrary
AEAD, and \cite{CR22} give \textsf{CTX}, which makes any nonce-based AE scheme
commit to its full context at the cost of a single hash over a short string and
with no ciphertext expansion. These transforms are inexpensive and broadly
applicable; \textsf{CTX} in particular is nearly free. \TW{} differs not by
undercutting their cost but by integration: it is committing as the mode itself,
with the same permutation and the same tag that already provide confidentiality
and integrity, so no separate commitment primitive, key-derivation step, or
auxiliary hash is invoked, and the strongest $s$-way full-input notion follows
from a single root tag collision bound (\Cref{sec:cmtds3}). More strongly,
\Cref{thm:tag-cmtds4} shows that the root tag alone is a compact commitment
string even when the ciphertext bodies used as openings differ, up to the
additional leaf tag birthday term. The same tag, being preimage-resistant,
additionally yields context discoverability against that fixed tag, which the
context-committing transforms are not known to provide. This places \TW{}
alongside the dedicated committing primitives---the compactly committing AEAD of
\cite{GLR17} and the encryptment of \cite{DGRW18}---which likewise build
commitment in rather than bolt it on, but which were designed for message
franking rather than as general nonce-based AEADs. The contrast with transformed
conventional schemes is sharpened by the attacks of \Cref{sec:related}: base
schemes such as GCM and OCB are not committing on their own \cite{CR22}, and a
transform is the only way to make them so.

Closest to \TW{} here is the recent work of \cite{DHMVV24}, which builds
nonce-based AE directly on the SHA-3 functions \textsf{SHAKE} and
\textsf{TurboSHAKE}---the latter using the same round-reduced \Keccak{}
permutation as \TW{}---and likewise achieves \textsf{CMT-4} natively, from the
collision resistance of the underlying function rather than from a transform. The
two designs make different structural choices: \cite{DHMVV24} targets
session-based communication, chaining a stream of messages through intermediate
tags and generalizing the keyed duplex into a \emph{duplex cipher}, whereas \TW{}
is a single two-level tree whose independent leaves expose run-time-selectable
parallelism at constant memory (\Cref{sec:related,sec:impl}). Because a session
chain processes each message serially, \TW{}'s principal advantage is throughput
on large messages (\Cref{sec:performance}): its parallel leaves fill vector units
a serial duplex leaves idle, at the cost of the session support \cite{DHMVV24}
provides and \TW{} does not. Both demonstrate
that committing AE needs no add-on transform once it is built on a
well-scrutinized \Keccak{} permutation.

\subsection{Limitations and Open Problems}
\label{sec:limitations}

\paragraph{Nonce misuse.}
\TW{}'s privacy analysis requires a nonce-respecting adversary, and the mode is
not misuse-resistant: a repeated nonce reuses leaf keystream and breaks
confidentiality, as for any online duplex stream cipher. Schemes such as
AES-GCM-SIV \cite{GLL17}, in the deterministic/misuse-resistant lineage of
\cite{RS06}, tolerate nonce repetition at the price of a second pass. A
misuse-resistant variant of \TW{} is left open.

\paragraph{Release of unverified plaintext.}
The analysis assumes that plaintext is released only after the tag verifies. \TW{}
is online and single-pass, so an implementation that streams decrypted leaf
output before checking the root tag operates in the release-of-unverified-plaintext
setting of \cite{ABLMMY14}, which our bounds do not cover; plaintext-awareness and
\textsf{INT-RUP} security for \TW{} are unstudied.

\paragraph{Bound tightness.}
Two terms are knowingly conservative. The flat-sponge differentiating loss
$\Lift_c(\Ntot) = \Ntot(\Ntot+1)/2^{c+1}$ bounds the indifferentiability
loss of \cite{BDPV08}, used here as a generic lift from the
random-oracle model. A direct analysis of the keyed duplex is tighter: the
modern keyed-duplex bounds synthesized by \cite{Men23}, building on the
full-state keyed-duplex line, bound a keyed duplex against an ideal permutation
without passing through indifferentiability, and recasting \TW{}'s leaf and root
duplexes in that framework is a promising route to a sharper capacity term. This
route is not turnkey, however: the keyed duplex bundles squeezing and absorption
into a single call, so its authenticated-encryption mode \textsf{MonkeySpongeWrap}
has the decryptor \emph{overwrite} the rate with each ciphertext block rather than
feed back the ciphertext it has just produced. On the final, padded block this
overwrite no longer matches the encryptor's XOR-absorbed padded plaintext outside
the message bits, so the decryptor cannot reconstruct the final block of duplex
output and recomputes the tag from a divergent state \cite[Alg.~8]{Men23}. \TW{}
sidesteps this by absorbing the ciphertext in both directions (\Cref{sec:tw128}),
so recasting it on the keyed duplex would first require repairing the mode. We
retain the indifferentiability lift for its modularity, as the resulting
$\Lift_c$ already meets the $128$-bit target (\Cref{sec:bound-discussion}). The
leaf tag collision term
$\Nleaf(\Nleaf-1)/2^{\tauleaf+1}$ is a pairwise union bound whose necessity for
the AE guarantee, as opposed to an artifact of the accounting, is not settled.

\paragraph{Beyond-birthday and leakage.}
All bounds are birthday-type on the $256$-bit capacity, delivering the
$128$-bit target but no more; a beyond-birthday analysis is open. The model is
also single-stage and leakage-free: side-channel resistance is addressed only at
the implementation level, through the constant-time properties of
\Cref{sec:impl}, and is not part of the provable-security claims.
